\section{String field theory, D-instantons, soft theorems and matrix theory}\label{sec:sft}

\subsection{String field theory}\label{ssec:sft}

Two claims of Sen about string field theory itself are in the catalogue. The first, with Zwiebach \cite{SenZwiebach94}, is that closed string field theory is background independent under infinitesimal deformations of the conformal field theory: the string field theories formulated around two nearby conformal backgrounds are related by a field redefinition. The second is the construction of superstring field theory \cite{Sen15offshell,Sen15onePI,Sen16BV} and the proof of its unitarity \cite{Sen16unitarity}: a gauge-invariant one-particle-irreducible effective action for heterotic and type II strings, free of the divergences of the picture-changing operators of earlier formulations through the use of a second, free string field of picture number $-3/2$; a Batalin--Vilkovisky master action; and the cutting rules, hence unitarity, of the resulting amplitudes.

\begin{verdict}[claims 30, 31]
Proved. Claim~30 is proved in \cite{SenZwiebach94} for marginal deformations to first order (and to second order in the companion paper), by constructing the field redefinition from the connection on the space of conformal field theories; its extension to finite deformations is a theorem only within the formal perturbation theory of the string field and is not a claim of the paper. Claim~31 is a construction and is proved in the papers that make it: the 1PI action is gauge invariant by the main identity of the vertices, the master action satisfies the master equation, and the Cutkosky rules are derived in \cite{Sen16unitarity} with the $i\epsilon$ prescription of \cite{Pius16}. The reader who wants the construction in one place has the review \cite{SenZwiebach24}. None of this depends on a conjecture; the open questions of superstring field theory are not about these statements.
\end{verdict}

\paragraphx{Corrected versions and infrared scope.} The published erratum to \cite{Pius16} and the corrected arXiv v3 must be used for the cutting argument; the correction is excluded from the list of unresolved errors. Section~4 of \cite{Sen16unitarity} explicitly notes the infrared problem for $D\le4$ and that the inclusive-cross-section construction is not carried out there. This is an acknowledged limitation, not a hidden failure of its perturbative cutting proof. A simple leading-soft cancellation and a counterexample to inferring exact unitarity from a formal series are given in Appendix~\ref{app:perturbative}. Appendix~\ref{app:infraredlimit} constructs the finite-energy limit in the explicit positive soft model and states the additional string-amplitude estimate needed for transfer. Appendix~\ref{app:softremainder} derives a uniform subtracted tree-emission bound from Sen's soft theorem and demonstrates why overlapping limits require additional estimates. None of these is a construction of the full infrared-finite string S-matrix.

\subsection{D-instantons}\label{ssec:instantons}

The contribution of D-instantons to string amplitudes has an overall normalisation that the world-sheet theory does not fix, because the open string path integral on the instanton has zero modes, bosonic and fermionic, whose measure is ambiguous. Sen \cite{Sen20fixing,Sen21normalisation,Sen21IIB} showed that open string field theory fixes it: the zero modes are treated by changing variables from the string field to the collective coordinates with a Jacobian computed in the field theory, the gauge zero mode of the $U(1)$ on the instanton is handled by a gauge fixing that the field theory provides, and the resulting normalisation is finite and unambiguous.

\begin{verdict}[claim 32]
Confirmed. In two-dimensional string theory the normalisation computed this way \cite{Sen20fixing,Sen21Dinst} agrees with the matrix model to the order compared, including the two-instanton contribution \cite{BRY23}; in type IIB string theory in ten dimensions the normalisation of the D-instanton correction to the four-graviton amplitude \cite{Sen21IIB} agrees with the coefficient predicted by $SL(2,\ZZ)$ invariance, the Green--Gutperle Eisenstein series. The claim that string field theory fixes the normalisation in every background is confirmed by these cases and not proved in general: the construction requires the string field theory of the instanton to be well defined, which it is in these examples, and a general statement needs a general treatment of the zero modes of the Ramond sector, which is not established by the examples discussed here.
\end{verdict}

\subsection{Soft theorems}\label{ssec:soft}

Laddha and Sen \cite{LaddhaSen19} proved the classical soft graviton theorem in dimensions $D>4$, the statement that the low-frequency gravitational radiation from a scattering process is determined by the momenta and angular momenta of the incoming and outgoing objects, including the sub-subleading term, directly from the classical equations \cite{Sen18subsub}; and Sahoo and Sen \cite{SahooSen19} found that in $D=4$ the soft expansion acquires logarithmic terms, $\ln\omega$ at the subleading order, whose coefficients are again universal, and conjectured the form of the corresponding quantum terms.

\begin{verdict}[claims 33, 34]
Claim~33 is proved: the classical theorem in $D>4$ is derived in \cite{LaddhaSen19}, and the classical logarithmic terms in $D=4$ in \cite{SahooSen19}, from the classical equations of motion with the long-range forces included. Claim~34 is confirmed: the quantum logarithmic terms were computed at one loop in \cite{SahooSen19} for scalar matter coupled to gravity and to electromagnetism and found to have the conjectured form, with the classical term and a quantum correction that depends only on the asymptotic momenta; the universality of the one-loop $\ln\omega$ terms was extended to spinning particles with the non-minimal couplings allowed by gauge and coordinate invariance by Krishna and Sahoo \cite{KrishnaSahoo23}, and the classical soft factors have been pushed to all orders in the electromagnetic case \cite{KKSS25}. The remaining conjecture is that the quantum $\ln\omega$ terms are one-loop exact, and the missing step is a proof that the two-loop and higher infrared divergences exponentiate in the way the conjecture requires; \cite{KrishnaSahoo23} states this as open, a complete higher-loop argument is not supplied here.
\end{verdict}

\paragraphx{A recent scope check.} The all-orders classical analysis of \cite{KKMSS26} distinguishes the Newtonian result from its proposed relativistic extension. It also finds that a separate all-orders conjecture of Alessio, Di Vecchia and Heissenberg, consistent with two-particle scattering, does not hold for general initial or final states with at least three particles. This is not a counterexample to Sen's earlier subleading quantum result. The extension in \cite{KKMSS26} is explicitly conjectural and is recorded as an open construction, not a wrong proof.

\subsection{Matrix theory}\label{ssec:matrix}

Sen \cite{Sen98matrix} argued, independently of Seiberg \cite{Seiberg97}, that the matrix theory description of M-theory in the infinite momentum frame follows from the compactification of M-theory on a light-like circle, taken as a limit of a small space-like circle, and extended it to tori $T^n$: the theory of $N$ D0-branes on $T^n$ in a limit in which the torus and the string scale shrink. For $n\le 3$ the limit is a maximally supersymmetric Yang--Mills theory in $n+1$ dimensions; for $n=4,5$ it is the $(2,0)$ theory on $T^5$ and the little string theory on $T^5$ respectively; and for $n\ge6$ the limit is, as Sen states, ``a theory without an independent definition'', the degrees of freedom of the D0-branes not decoupling from gravity. The two statements H1 and H2 of the catalogue are these, and they are recorded without a verdict: the first is the Seiberg--Sen argument, which is a derivation within the assumptions of the infinite momentum frame and has no independent test for $n\ge 4$; the second is a statement about a theory that is not defined, and cannot be tested in any regime. We mention them because the paper is cited as a derivation of matrix theory and the reader should know that its content for $n\ge6$ is a limit of a non-decoupled system, as its author says.
